% N3 --- Superposicao: 10 questoes com topologias distintas
% Revisao visual: ramos mais espaçados, rotulos afastados e sem diagonais
% desnecessarias ou sobreposicao entre componentes.
\blocoaprofundamento

\begin{questao}[3]{}
A fonte dependente deve permanecer ativa em todas as parcelas. Determine $V_A$ por superposição das duas fontes independentes.
\circuitoq{\begin{circuitikz}[american,scale=.74,transform shape,line width=.8pt]
\coordinate (g) at (4.5,0); \coordinate (a) at (4.5,4.4);
\draw (0,0) to[\fonteV,l^={$\SI{12}{\volt}$}] (0,4.4)
      to[R,l={$\SI{4}{\ohm}$}] (a);
\draw (a) to[R,l_={$\SI{8}{\ohm}$}] (g)--(0,0);
\draw (8.0,0) to[isource,l_={$\SI{1}{\ampere}$}] (8.0,4.4)--(a);
\draw (11.6,0) to[cisource,l_={$0{,}0625V_A$}] (11.6,4.4)--(8.0,4.4);
\draw (11.6,0)--(8.0,0)--(g);
\fill (a) circle (1.1pt);
\node[Vcol,fill=white,inner sep=1.6pt] at ($(a)+(0,.68)$) {$V_A$};
\node[ground] at(g){};
\end{circuitikz}}
\resposta{$V_A'=\SI{9.6}{\volt}$, $V_A''=\SI{3.2}{\volt}$, $V_A=\SI{12.8}{\volt}$}
\resolucao{A condutância líquida é $1/4+1/8-0{,}0625=0{,}3125$ S. A fonte de 12 V injeta equivalente de 3 A, dando $9{,}6$ V; a fonte de 1 A dá $3{,}2$ V. A dependente permanece nas duas análises.}
\end{questao}

\begin{questao}[3]{}
Há três fontes independentes atuando em dois nós. Determine $V_1$ e $V_2$ por superposição.
\circuitoq{\begin{circuitikz}[american,scale=.62,transform shape,line width=.8pt]
\coordinate (g) at (6.2,0); \coordinate (n1) at (4.1,4.8); \coordinate (n2) at (8.5,4.8);
\draw (0,0) to[\fonteV,l^={$\SI{18}{\volt}$}] (0,4.8)
      to[R,l={$\SI{6}{\ohm}$}] (n1);
\draw (2.0,0) to[isource,l^={$\SI{2}{\ampere}$}] (2.0,4.8)--(n1);
\draw (n1) to[R,l_={$\SI{9}{\ohm}$}] (4.1,0)--(g);
\draw (n1) to[R,l={$\SI{3}{\ohm}$}] (n2);
\draw (n2) to[R,l={$\SI{6}{\ohm}$}] (8.5,0)--(g);
\draw (12.8,0) to[\fonteV,l_={$\SI{9}{\volt}$}] (12.8,4.8)
      to[R,l_={$\SI{3}{\ohm}$}] (n2);
\draw (0,0)--(2.0,0)--(g)--(12.8,0);
\fill (n1) circle (1.1pt); \fill (n2) circle (1.1pt);
\node[Vcol,fill=white,inner sep=1.6pt] at ($(n1)+(0,.72)$) {$V_1$};
\node[Vcol,fill=white,inner sep=1.6pt] at ($(n2)+(0,.72)$) {$V_2$};
\node[ground] at(g){};
\end{circuitikz}}
\resposta{$V_1\approx\SI{12.98}{\volt}$; $V_2\approx\SI{8.79}{\volt}$}
\resolucao{Resolvem-se três vezes as mesmas duas equações nodais, uma por fonte. A soma das parcelas fornece $V_1=558/43\approx12{,}98$ V e $V_2=378/43\approx8{,}79$ V.}
\end{questao}

\begin{questao}[3]{}
Duas fontes reais alimentam uma carga. Determine a tensão $U$ e identifique a parcela de cada f.e.m.
\circuitoq{\begin{circuitikz}[american,scale=.74,transform shape,line width=.8pt]
\coordinate (g) at (5.3,0); \coordinate (a) at (5.3,4.4);
\draw (0,0) to[\fonteV,l^={$\SI{18}{\volt}$}] (0,4.4)
      to[R,l={$r_1=\SI{3}{\ohm}$}] (a);
\draw (11.0,0) to[\fonteV,l_={$\SI{9}{\volt}$}] (11.0,4.4)
      to[R,l_={$r_2=\SI{6}{\ohm}$}] (a);
\draw (a) to[R,l_={$R_L=\SI{9}{\ohm}$}] (g);
\draw (0,0)--(g)--(11.0,0);
\fill (a) circle (1.1pt);
\node[Vcol,fill=white,inner sep=1.6pt] at ($(a)+(0,.68)$) {$U$};
\node[ground] at(g){};
\end{circuitikz}}
\resposta{$U_1\approx\SI{9.82}{\volt}$; $U_2\approx\SI{2.45}{\volt}$; $U\approx\SI{12.27}{\volt}$}
\resolucao{A condutância total é $1/3+1/6+1/9=11/18$ S. As parcelas são $(18/3)/(11/18)=108/11$ V e $(9/6)/(11/18)=27/11$ V. Soma $135/11\approx12{,}27$ V.}
\end{questao}

\begin{questao}[3]{}
A fonte de tensão controlada impõe $V_2=1{,}5V_1$. Determine $V_1$ por superposição das fontes independentes.
\circuitoq{\begin{circuitikz}[american,scale=.72,transform shape,line width=.8pt]
\coordinate (g) at (3.0,0); \coordinate (n1) at (3.0,4.5); \coordinate (n2) at (8.2,4.5);
\draw (0,0) to[isource,l^={$\SI{2}{\ampere}$}] (0,4.5)--(n1);
\draw (n1) to[R,l_={$\SI{6}{\ohm}$}] (g);
\draw (n1) to[R,l={$\SI{12}{\ohm}$}] (n2);
\draw (n2) to[cvsource,l_={$1{,}5V_1$}] (8.2,0)--(g)--(0,0);
\draw (11.2,0) to[isource,l_={$\SI{0.5}{\ampere}$}] (11.2,4.5)--(n2);
\draw (11.2,0)--(8.2,0);
\fill (n1) circle (1.1pt); \fill (n2) circle (1.1pt);
\node[Vcol,fill=white,inner sep=1.6pt] at ($(n1)+(0,.72)$) {$V_1$};
\node[Vcol,fill=white,inner sep=1.6pt] at ($(n2)+(0,.72)$) {$V_2$};
\node[ground] at(g){};
\end{circuitikz}}
\resposta{A fonte dependente permanece ativa em ambas as análises; $V_1=\SI{16}{\volt}$}
\resolucao{A superposição é feita somente sobre as duas fontes independentes; a fonte controlada não é zerada. Montando as equações em cada parcela e somando, obtém-se $V_1=16$ V e $V_2=24$ V.}
\end{questao}

\begin{questao}[3]{}
Mostre quantitativamente por que não se deve somar potências parciais no resistor.
\circuitoq{\begin{circuitikz}[american,scale=.90,transform shape,line width=.8pt]
\draw (0,0) to[\fonteV,l^={$\SI{16}{\volt}$}] (0,3.8)
      to[R,l={$R=\SI{4}{\ohm}$},i>^={$I$}] (5.6,3.8)
      to[\fonteV,l_={$\SI{8}{\volt}$}] (5.6,0)--(0,0);
\end{circuitikz}}
\resposta{$I_1=\SI{4}{\ampere}$, $I_2=\SI{2}{\ampere}$; $P=\SI{144}{\watt}$, enquanto $P_1+P_2=\SI{80}{\watt}$}
\resolucao{As correntes somam: $I=4+2=6$ A. Portanto $P=RI^2=4(36)=144$ W. Somar $4(4^2)+4(2^2)=80$ W ignora o termo cruzado $2RI_1I_2=64$ W.}
\end{questao}

\begin{questao}[3]{}
Use o circuito para explicar por que a superposição falha quando o estado do diodo muda entre parcelas.
\circuitoq{\begin{circuitikz}[american,scale=.74,transform shape,line width=.8pt]
\coordinate (g) at (5.0,0); \coordinate (a) at (5.0,4.4);
\draw (0,0) to[\fonteV,l^={$\SI{6}{\volt}$}] (0,4.4)
      to[R,l={$\SI{1}{\kilo\ohm}$}] (a);
\draw (10.4,0) to[\fonteV,l_={$\SI{4}{\volt}$},invert] (10.4,4.4)
      to[R,l_={$\SI{2}{\kilo\ohm}$}] (a);
\draw (a) to[D,l_={$D$}] (g);
\draw (0,0)--(g)--(10.4,0);
\fill (a) circle (1.1pt); \node[ground] at(g){};
\end{circuitikz}}
\resposta{Não se pode superpor as respostas obtidas com estados diferentes do diodo}
\resolucao{O diodo altera a própria topologia equivalente ao conduzir ou bloquear. Assim, cada parcela pode corresponder a uma rede diferente; a linearidade exigida pelo teorema deixa de existir.}
\end{questao}

\begin{questao}[3]{}
No circuito linear, determine o novo $V_A$ se a fonte de $\SI{4}{\volt}$ for duplicada, sem refazer toda a análise.
\circuitoq{\begin{circuitikz}[american,scale=.76,transform shape,line width=.8pt]
\coordinate (g) at (5.2,0); \coordinate (a) at (5.2,4.3);
\draw (0,0) to[\fonteV,l^={$\SI{10}{\volt}$}] (0,4.3)
      to[R,l={$\SI{5}{\ohm}$}] (a);
\draw (10.8,0) to[\fonteV,l_={$\SI{4}{\volt}$}] (10.8,4.3)
      to[R,l_={$\SI{10}{\ohm}$}] (a);
\draw (a) to[R,l_={$\SI{10}{\ohm}$}] (g);
\draw (0,0)--(g)--(10.8,0);
\fill (a) circle (1.1pt);
\node[Vcol,fill=white,inner sep=1.6pt] at ($(a)+(0,.68)$) {$V_A$};
\node[ground] at(g){};
\end{circuitikz}}
\resposta{Originalmente $V_A=\SI{6}{\volt}$; após duplicar a segunda fonte, $V_A=\SI{7}{\volt}$}
\resolucao{As parcelas originais são $5$ V e $1$ V. Duplicar apenas a segunda fonte duplica somente sua parcela: $1\to2$ V. Logo o novo total é $5+2=7$ V.}
\end{questao}

\begin{questao}[3]{}
A corrente total no galvanômetro é nula. A fonte $E_1$, sozinha, produz $+\SI{35}{\milli\ampere}$. Determine a contribuição de $E_2$ e interprete.
\circuitoq{\begin{circuitikz}[american,scale=.62,transform shape,line width=.8pt]
\coordinate (l) at (0,2.5); \coordinate (r) at (12.0,2.5);
\coordinate (u) at (6.0,6.2); \coordinate (d) at (6.0,-1.2);
\draw (l) to[\fonteV,l^={$E_1$}] (0,-1.2)--(d);
\draw (r) to[\fonteV,l_={$E_2$}] (12.0,-1.2)--(d);
\draw (l) to[R,l={$R_1$}] (u) to[R,l={$R_3$}] (r);
\draw (l) to[R,l_={$R_2$}] (d) to[R,l_={$R_4$}] (r);
\draw (u) to[R,l={$R_g$},i>^={$I_g$}] (d);
\node[Vcol,fill=white,inner sep=1.5pt] at ($(u)+(0,.55)$) {$a$};
\node[Vcol,fill=white,inner sep=1.5pt] at ($(d)+(0,-.55)$) {$b$};
\node[ground] at(d){};
\end{circuitikz}}
\resposta{$I_g^{(E_2)}=\SI{-35}{\milli\ampere}$}
\resolucao{Como $I_g=I_g^{(E_1)}+I_g^{(E_2)}=0$, a segunda parcela deve ser $-35$ mA. Isso é cancelamento entre fontes, não necessariamente equilíbrio resistivo da ponte.}
\end{questao}

\begin{questao}[3]{}
Um sinal e uma interferência entram por ramos diferentes. Determine as parcelas de saída e a razão sinal-interferência.
\circuitoq{\begin{circuitikz}[american,scale=.72,transform shape,line width=.8pt]
\coordinate (g) at (5.3,0); \coordinate (a) at (5.3,4.5);
\draw (0,0) to[\fonteV,l^={$V_s=\SI{0.2}{\volt}$}] (0,4.5)
      to[R,l={$\SI{1}{\kilo\ohm}$}] (a);
\draw (11.0,0) to[\fonteV,l_={$V_n=\SI{1}{\volt}$}] (11.0,4.5)
      to[R,l_={$\SI{20}{\kilo\ohm}$}] (a);
\draw (a) to[R,l_={$R_L=\SI{2}{\kilo\ohm}$}] (g);
\draw (0,0)--(g)--(11.0,0);
\fill (a) circle (1.1pt);
\node[Vcol,fill=white,inner sep=1.6pt] at ($(a)+(0,.70)$) {$V_o$};
\node[ground] at(g){};
\end{circuitikz}}
\resposta{$V_{o,s}\approx\SI{129}{\milli\volt}$; $V_{o,n}\approx\SI{32.3}{\milli\volt}$; razão $\approx4$ ($\approx\SI{12}{\decibel}$)}
\resolucao{A condutância total é $1/1k+1/20k+1/2k=1{,}55$ mS. As injeções são $0{,}2$ mA e $0{,}05$ mA, gerando aproximadamente $129$ mV e $32{,}3$ mV.}
\end{questao}

\begin{questao}[3]{}
Compare o custo de superposição e análise nodal para a rede de dois nós e quatro fontes abaixo.
\circuitoq{\begin{circuitikz}[american,scale=.56,transform shape,line width=.8pt]
\coordinate (g) at (7.0,0); \coordinate (n1) at (4.4,5.0); \coordinate (n2) at (9.4,5.0);
\draw (0,0) to[\fonteV,l^={$E_1$}] (0,5.0) to[R,l={$R_1$}] (n1);
\draw (2.1,0) to[isource,l^={$I_2$}] (2.1,5.0)--(n1);
\draw (n1) to[R,l_={$R_2$}] (4.4,0)--(g);
\draw (n1) to[R,l={$R_{12}$}] (n2);
\draw (n2) to[R,l={$R_3$}] (9.4,0)--(g);
\draw (11.8,0) to[isource,l_={$I_4$}] (11.8,5.0)--(n2);
\draw (14.0,0) to[\fonteV,l_={$E_3$}] (14.0,5.0) to[R,l_={$R_4$}] (n2);
\draw (0,0)--(2.1,0)--(g)--(11.8,0)--(14.0,0);
\fill (n1) circle (1.1pt); \fill (n2) circle (1.1pt);
\node[Vcol,fill=white,inner sep=1.5pt] at ($(n1)+(0,.72)$) {$V_1$};
\node[Vcol,fill=white,inner sep=1.5pt] at ($(n2)+(0,.72)$) {$V_2$};
\node[ground] at(g){};
\end{circuitikz}}
\resposta{Superposição: quatro sistemas $2\times2$; nodal direto: um sistema $2\times2$}
\resolucao{Para um único resultado numérico, nodal é mais econômico. Superposição vale o custo adicional quando é necessário conhecer a contribuição individual de cada fonte.}
\end{questao}